\documentclass[11pt,a4paper]{article}

\usepackage[utf8]{inputenc}
\usepackage[T1]{fontenc}
\usepackage{amsmath,amssymb,bm}
\usepackage{booktabs}
\usepackage[margin=2.5cm]{geometry}
\usepackage{hyperref}
\usepackage{enumitem}

\hypersetup{
    colorlinks=true,
    linkcolor=blue,
    citecolor=blue,
    urlcolor=blue
}

\newcommand{\mb}[1]{\mathbf{#1}}
\newcommand{\phat}{\hat{\mb{p}}}
\newcommand{\zhat}{\hat{z}}
\newcommand{\Heff}{\mb{H}_{\text{eff}}}
\newcommand{\Hex}{\mb{H}_{\text{ex}}}
\newcommand{\Hdmi}{\mb{H}_{\text{DMI}}}
\newcommand{\Hanis}{\mb{H}_{\text{anis}}}
\newcommand{\Hdemag}{\mb{H}_{\text{demag}}}
\newcommand{\Hrkky}{\mb{H}_{\text{RKKY}}}
\newcommand{\Hext}{\mb{H}_{\text{ext}}}
\newcommand{\tCo}{t_{\text{Co}}}
\newcommand{\Keff}{K_{\text{eff}}}
\newcommand{\Aex}{A_{\text{ex}}}
\newcommand{\Ms}{M_s}

\title{Theoretical Background:\\
SAF Skyrmion Simulator}
\author{Rui Barreira}
\date{April 2026}

\begin{document}
\maketitle

%======================================================================
\section{Magnetic Skyrmions in Synthetic Antiferromagnets}
%======================================================================

Magnetic skyrmions are topologically non-trivial chiral spin textures
with a topological charge $S = \pm 1$. In a skyrmion, the
magnetization rotates smoothly from one out-of-plane direction at the
core to the opposite direction in the surrounding domain.

A synthetic antiferromagnet (SAF) consists of two ferromagnetic layers
(Co) separated by a non-magnetic spacer (Ru), coupled
antiferromagnetically via the Ruderman--Kittel--Kasuya--Yosida (RKKY)
interaction. In a SAF, skyrmions in the two layers have opposite core
polarities and opposite topological charges, leading to a cancellation
of the Magnus force and a vanishing skyrmion Hall
effect~\cite{pham2024}.

The SAF stack modeled here is Pt/Co/Ru/Pt/Co/Ru, following the
experimental system in~\cite{pham2024}.

%----------------------------------------------------------------------
\subsection{Notation}
\label{sec:notation}

\begin{itemize}[leftmargin=2em]
    \item $\mb{M}(\mb{r}, t)$ --- magnetization.
    \item $\Ms$ --- saturation magnetization at the operating
        temperature, held fixed.
    \item $\mb{m} = \mb{M}/\Ms$ --- unit magnetization vector,
        $|\mb{m}| = 1$; evolved by the LLGS equation.
    \item $\Heff, \Hex, \Hdmi, \Hanis, \Hdemag, \Hrkky, \Hext$ ---
        effective fields entering the LLGS equation
        $d\mb{m}/dt = -\gamma\,\mb{m}\times\Heff$. The
        $\mb{B}=\mu_0\mb{H}$ factor is absorbed into the prefactors
        (B-field units; see Sec.~\ref{sec:units}).
    \item $J_{\text{RKKY}}$ --- RKKY interlayer coupling constant;
        surface energy density (J/m$^2$) in Section~2.6, volume
        energy density (J/m$^3$) in Sec.~\ref{sec:variational},
        related by $J_{\text{RKKY}}^{(\text{vol})}
        = J_{\text{RKKY}}^{(\text{surf})}/\tCo$.
    \item $H_{\text{RKKY}}$ --- RKKY field magnitude:
        $H_{\text{RKKY}} = J_{\text{RKKY}}^{(\text{vol})}/\Ms
        = J_{\text{RKKY}}^{(\text{surf})}/(\Ms\,\tCo)$.
    \item $\mb{J} = J\,\hat{J}$ --- current density;
        $J$ is the magnitude, $\hat{J}$ the unit vector along the
        current flow.
    \item $\phat$ --- spin polarization direction,
        $\phat = \zhat \times \hat{J}$.
    \item $\zhat$ --- film-normal unit vector (out-of-plane).
    \item $H_{\text{DL}} = \chi_{\text{DL}} J$,
        $H_{\text{FL}} = \chi_{\text{FL}} J$ --- damping-like and
        field-like SOT effective fields.
    \item $\Keff = K - \mu_0 \Ms^2 / 2$ --- effective uniaxial
        anisotropy with thin-film shape correction.
    \item $V_{\text{cell}} = \tCo\,a^2$ --- per-site simulation
        volume.
    \item Energy integrals over a layer:
        $\int \cdots\,dV = \tCo \int \cdots\,dS$ for thin uniform
        layers.
\end{itemize}

%----------------------------------------------------------------------
\subsection{Units}
\label{sec:units}

All quantities are in SI. Effective fields are carried in Tesla
(B-field units), absorbing the $\mu_0$ factor from
$\mb{B} = \mu_0\mb{H}$.

\begin{table}[h!]
\centering
\begin{tabular}{@{}lll@{}}
\toprule
Symbol & Quantity & Unit \\
\midrule
$\mb{M}, \Ms$ & magnetization, saturation magnetization & A/m \\
$\mb{m}$ & unit magnetization vector & --- \\
$d\mb{m}/dt$ & magnetization rate of change & s$^{-1}$ \\
$\Heff, \Hex, \Hdmi, \Hanis$ & effective fields & T \\
$\Hdemag, \Hrkky, \Hext, \mb{B}_{\text{ext}}$ & effective fields & T \\
$H_{\text{DL}}, H_{\text{FL}}$ & SOT effective fields & T \\
$\Aex$ & exchange stiffness & J/m \\
$D$ & DMI constant & J/m$^2$ \\
$K, K_{\text{top}}, K_{\text{bot}}, \Keff$ & uniaxial anisotropy
    & J/m$^3$ \\
$\alpha$ & Gilbert damping & --- \\
$\gamma$ & gyromagnetic ratio & rad/(s$\cdot$T) \\
$\mu_0$ & vacuum permeability & T$\cdot$m/A \\
$J_{\text{RKKY}}^{(\text{surf})}$ & RKKY surface energy density
    & J/m$^2$ \\
$J_{\text{RKKY}}^{(\text{vol})}$ & RKKY volume energy density
    & J/m$^3$ \\
$H_{\text{RKKY}}$ & RKKY field magnitude & T \\
$\mb{J}, J$ & current density, magnitude & A/m$^2$ \\
$\hat{J}, \phat, \zhat$ & unit vectors & --- \\
$\chi_{\text{DL}}, \chi_{\text{FL}}$ & SOT coefficients
    & T$\cdot$A$^{-1}\cdot$m$^2$ \\
$a$ & lattice constant & m \\
$\tCo, d_{\text{Ru}}$ & layer / spacer thicknesses & m \\
$\Delta$ & domain wall width & m \\
$R$ & skyrmion radius & m \\
$V_{\text{cell}} = \tCo\,a^2$ & per-site simulation volume & m$^3$ \\
$n_x, n_y$ & lattice sites & --- \\
$t, \Delta t$ & time, time step & s \\
$\mathcal{H}$ & Hamiltonian (energy) & J \\
$Q$ & topological charge & --- \\
$\mb{G}$ & gyrovector & --- \\
\bottomrule
\end{tabular}
\end{table}

%======================================================================
\section{Hamiltonian}
%======================================================================

The total magnetic energy of the SAF system is the sum of
contributions from both layers
($\ell = \text{top}, \text{bot}$) plus the interlayer coupling:
%
\begin{equation}
\mathcal{H} = \sum_{\ell} \left(
    \mathcal{H}_{\text{ex}}^{(\ell)}
    + \mathcal{H}_{\text{DMI}}^{(\ell)}
    + \mathcal{H}_{\text{anis}}^{(\ell)}
    + \mathcal{H}_{\text{demag}}^{(\ell)}
    + \mathcal{H}_{\text{Z}}^{(\ell)}
\right)
+ \mathcal{H}_{\text{RKKY}}
\end{equation}

The effective field entering the LLGS equation is obtained from the
variational derivative of the energy with respect to the
magnetization:
%
\begin{equation}
\Heff = -\frac{1}{\Ms}
    \frac{\delta \mathcal{H}}{\delta \mb{m}}
\end{equation}

%----------------------------------------------------------------------
\subsection{Exchange Energy}

The Heisenberg exchange favors parallel alignment of neighboring
spins:
%
\begin{equation}
\mathcal{H}_{\text{ex}} = \Aex \int \left[
    (\nabla m_x)^2 + (\nabla m_y)^2 + (\nabla m_z)^2
\right] dV
\end{equation}

On the discrete 2D lattice with lattice constant $a$ and film
thickness $\tCo$:
%
\begin{equation}
\mathcal{H}_{\text{ex}} = -\Aex \, \tCo
    \sum_{\langle i,j \rangle} \mb{m}_i \cdot \mb{m}_j
\end{equation}
%
where the sum runs over nearest-neighbor pairs.

%----------------------------------------------------------------------
\subsection{Dzyaloshinskii--Moriya Interaction Energy}

The interfacial (N\'{e}el-type) DMI at the Pt/Co interface introduces
chirality:
%
\begin{equation}
\mathcal{H}_{\text{DMI}} = D \, \tCo \int \left[
    m_z \, (\nabla \cdot \mb{m})
    - (\mb{m} \cdot \nabla) m_z
\right] dS
\end{equation}

On the discrete lattice this becomes a sum over nearest-neighbor bonds
with bond-direction-dependent cross products:
%
\begin{equation}
\mathcal{H}_{\text{DMI}} = -D \, \tCo
    \sum_{\langle i,j \rangle}
    (\hat{r}_{ij} \times \zhat) \cdot
    (\mb{m}_i \times \mb{m}_j)
\end{equation}
%
where $\hat{r}_{ij}$ is the unit vector from site $i$ to site $j$.

%----------------------------------------------------------------------
\subsection{Anisotropy Energy}

Uniaxial perpendicular magnetic anisotropy (PMA) along the film
normal:
%
\begin{equation}
\mathcal{H}_{\text{anis}} = -K \, \tCo \int m_z^2 \, dS
\end{equation}

On the lattice:
%
\begin{equation}
\mathcal{H}_{\text{anis}} = -K \, \tCo \, a^2 \sum_i m_{z,i}^2
\end{equation}

PMA in Co/Pt and Co/Ru/Pt stacks originates from interfacial
spin-orbit coupling at the heavy-metal interface (3d--5d
hybridization between Co and Pt). The effect scales with the
interface area, so thinner Co (here $\tCo = 1.3$~nm) has a larger
interface-to-volume ratio and stronger PMA per unit volume.

%----------------------------------------------------------------------
\subsection{Demagnetization Energy}

In the thin-film limit ($\tCo \ll$ lateral dimensions), the
magnetostatic energy reduces to a local shape anisotropy:
%
\begin{equation}
\mathcal{H}_{\text{demag}} = \frac{\mu_0 \Ms^2}{2} \, \tCo
    \int m_z^2 \, dS
\end{equation}

This opposes out-of-plane magnetization.  Since both the anisotropy
and demagnetization energies are proportional to $m_z^2$, they merge
into a single effective anisotropy:
%
\begin{equation}
\mathcal{H}_{\text{anis}} + \mathcal{H}_{\text{demag}} =
    -\underbrace{\left(K - \frac{\mu_0 \Ms^2}{2}\right)}_{\Keff}
    \, \tCo \int m_z^2 \, dS
\end{equation}
%
Three regimes arise depending on the sign of $\Keff$:
%
\begin{itemize}[nosep]
\item $K > \mu_0 \Ms^2/2$: $\Keff > 0$
    --- anisotropy wins, magnetization points out-of-plane.
\item $K < \mu_0 \Ms^2/2$: $\Keff < 0$
    --- demagnetization wins, magnetization lies in-plane.
\item $K \approx \mu_0 \Ms^2/2$: $\Keff \approx 0$
    --- spin reorientation transition (SRT); domain walls and
    skyrmions cost very little energy to form.
\end{itemize}
%
For the Co/Pt system in~\cite{pham2024},
$K \approx 1.3 \times 10^6$~J/m$^3$ and
$\mu_0 \Ms^2/2 \approx 1.286 \times 10^6$~J/m$^3$, giving
$\Keff \approx 10^4$~J/m$^3$ --- barely positive, right at the SRT.
This is why the system can host skyrmions.

%----------------------------------------------------------------------
\subsubsection*{Interface PMA vs shape anisotropy: the spin
reorientation transition}

In a thin ferromagnetic film, two anisotropies compete:
\begin{enumerate}[leftmargin=2em, nosep]
    \item \textbf{Interfacial PMA} ($K_s$, J/m$^2$) --- created by
        spin-orbit hybridization at heavy-metal interfaces (Pt, Ir,
        W, Pd, Au). Favors out-of-plane.
    \item \textbf{Shape anisotropy}
        ($\tfrac{1}{2}\mu_0 \Ms^2$, J/m$^3$) --- from the demag
        field of a thin slab uniformly magnetized along $\zhat$.
        Favors in-plane.
\end{enumerate}

The per-volume effective anisotropy of a film of thickness $t$ with
two equivalent interfaces is
%
\begin{equation}\label{eq:Keff_full}
\Keff = K_{\text{vol}} + \frac{2 K_s}{t}
        - \frac{1}{2}\mu_0 \Ms^2,
\end{equation}
%
where:
\begin{itemize}[leftmargin=2em, nosep]
    \item $K_s/t$ scales as $1/t$ --- interfaces dominate as $t$
        shrinks.
    \item $\tfrac{1}{2}\mu_0 \Ms^2$ is independent of thickness ---
        pure volume term.
\end{itemize}

Sign convention:
\begin{itemize}[leftmargin=2em, nosep]
    \item $\Keff > 0 \;\Longrightarrow\;$ out-of-plane easy axis
        (PMA).
    \item $\Keff < 0 \;\Longrightarrow\;$ in-plane easy axis.
\end{itemize}

There is a critical thickness $t_c$ where $\Keff$ changes sign:
%
\begin{equation}\label{eq:t_c}
t_c = \frac{2 K_s}{\tfrac{1}{2}\mu_0 \Ms^2 - K_{\text{vol}}}.
\end{equation}
%
For Co/Pt: $K_s \sim 0.5$--$1$~mJ/m$^2$,
$\mu_0 \Ms^2/2 \sim 1.3 \times 10^6$~J/m$^3$, $K_{\text{vol}}$ tiny
$\;\Longrightarrow\; t_c \sim 1$~nm.
\begin{itemize}[leftmargin=2em, nosep]
    \item $t < t_c$: PMA wins, magnetization out-of-plane.
    \item $t > t_c$: shape anisotropy wins, magnetization in-plane.
\end{itemize}

The simulator's Co layer is $\tCo = 1.3$~nm --- \textbf{right at
the SRT, deliberately}. That is why
$\Keff = K_{\text{top}} - \mu_0 \Ms^2/2 \approx 8\times10^3$~J/m$^3$
is a small difference of two large numbers ($\sim 10^6$ each), and
why 7\% uncertainty in $K$ propagates to $\sim$100\% uncertainty in
$\Keff$ (see Implementation Notes,
Sec.~\ref{sec:implementation_notes}).

%----------------------------------------------------------------------
\subsection{Zeeman Energy}

Interaction with an external field $\mb{B}_{\text{ext}}$:
%
\begin{equation}
\mathcal{H}_{\text{Z}} = -\Ms \, \tCo \int
    \mb{m} \cdot \mb{B}_{\text{ext}} \, dS
\end{equation}

%----------------------------------------------------------------------
\subsection{RKKY Interlayer Coupling Energy}

The antiferromagnetic RKKY exchange between the two Co layers across
the Ru spacer:
%
\begin{equation}
\mathcal{H}_{\text{RKKY}} = J_{\text{RKKY}} \int
    \mb{m}_{\text{top}} \cdot \mb{m}_{\text{bot}} \, dS
\end{equation}
%
where $J_{\text{RKKY}} > 0$ for antiferromagnetic coupling (surface
energy density, J/m$^2$). In the Tesla convention adopted here, the
coupling field magnitude is
$H_{\text{RKKY}} = J_{\text{RKKY}} / (\Ms \, \tCo)$ (cf.\ Notation
block).

%----------------------------------------------------------------------
\subsection{Total Energy on the Discrete Lattice}

Combining all contributions, the total energy per layer on the 2D
lattice is:
%
\begin{equation}
\mathcal{H}^{(\ell)} = \tCo \, a^2 \sum_i \left[
    -\frac{\Aex}{a^2} \sum_{\delta}
        \mb{m}_i \cdot \mb{m}_{i+\delta}
    - \frac{D}{a} \sum_{\delta}
        (\hat{\delta} \times \zhat) \cdot
        (\mb{m}_i \times \mb{m}_{i+\delta})
    - \Keff \, m_{z,i}^2
    - \Ms \, \mb{m}_i \cdot \mb{B}_{\text{ext}}
\right]
\end{equation}
%
where $\delta$ runs over the four nearest neighbors
$\{\pm\hat{x}, \pm\hat{y}\}$, and the interlayer RKKY term adds
$J_{\text{RKKY}} \, a^2 \sum_i
    \mb{m}_{\text{top},i} \cdot \mb{m}_{\text{bot},i}$.

%======================================================================
\section{From Hamiltonian to Effective Field}
\label{sec:variational}
%======================================================================

The Hamiltonian (Section~2) defines \emph{what} the system wants to
minimize.  The LLGS equation (Section~4) describes \emph{how} the
magnetization evolves in time.  The bridge between the two is the
\textbf{effective field}, obtained from the variational derivative of
the energy with respect to the magnetization.

%----------------------------------------------------------------------
\subsection{Variational Derivative}

Consider a small perturbation
$\mb{m} \to \mb{m} + \delta\mb{m}$ (subject to the constraint
$|\mb{m}| = 1$, so $\mb{m} \cdot \delta\mb{m} = 0$).  The
corresponding change in energy is:
%
\begin{equation}
\delta\mathcal{H} = \int
    \frac{\delta\mathcal{H}}{\delta\mb{m}}
    \cdot \delta\mb{m} \; dV
\end{equation}
%
The effective magnetic field is defined such that:
%
\begin{equation}\label{eq:heff_def}
\boxed{\;
\Heff = -\frac{1}{\Ms}
    \frac{\delta\mathcal{H}}{\delta\mb{m}}
\;}
\end{equation}
%
so that $\delta\mathcal{H} = -\Ms \int \Heff \cdot
\delta\mb{m} \; dV$.  The factor $1/\Ms$ gives $\Heff$ the
dimensions of magnetic field (Tesla in our convention).

%----------------------------------------------------------------------
\subsection{Derivation of the Exchange Field}

Starting from the exchange energy:
%
\begin{equation}
\mathcal{H}_{\text{ex}} = \Aex \int
    |\nabla\mb{m}|^2 \; dV
\end{equation}
%
Perturbing $\mb{m} \to \mb{m} + \delta\mb{m}$:
%
\begin{equation}
\delta\mathcal{H}_{\text{ex}} = 2\Aex \int
    \nabla\mb{m} : \nabla(\delta\mb{m}) \; dV
\end{equation}
%
where the colon denotes the tensor double contraction
$\nabla\mb{m} : \nabla(\delta\mb{m})
= \sum_\alpha \nabla m_\alpha \cdot \nabla(\delta m_\alpha)$.
Integrating by parts (boundary terms vanish for periodic boundary
conditions):
%
\begin{equation}
\delta\mathcal{H}_{\text{ex}} =
    -2\Aex \int \nabla^2\mb{m} \cdot \delta\mb{m} \; dV
\end{equation}
%
Reading off the variational derivative and applying~\eqref{eq:heff_def}:
%
\begin{equation}
\frac{\delta\mathcal{H}_{\text{ex}}}{\delta\mb{m}}
    = -2\Aex \, \nabla^2\mb{m}
\qquad\Longrightarrow\qquad
\Hex = \frac{2\Aex}{\Ms} \nabla^2\mb{m}
\end{equation}

%----------------------------------------------------------------------
\subsection{Derivation of the DMI Field}

From the interfacial DMI energy:
%
\begin{equation}
\mathcal{H}_{\text{DMI}} = D \int \left[
    m_z (\nabla \cdot \mb{m})
    - (\mb{m} \cdot \nabla) m_z
\right] dV
\end{equation}
%
Expanding each term under the variation
$\mb{m} \to \mb{m} + \delta\mb{m}$ and collecting contributions
from $\delta m_x$, $\delta m_y$, and $\delta m_z$ separately (using
integration by parts where needed):
%
\begin{equation}
\frac{\delta\mathcal{H}_{\text{DMI}}}{\delta m_x}
    = -2D\,\frac{\partial m_z}{\partial x}, \quad
\frac{\delta\mathcal{H}_{\text{DMI}}}{\delta m_y}
    = -2D\,\frac{\partial m_z}{\partial y}, \quad
\frac{\delta\mathcal{H}_{\text{DMI}}}{\delta m_z}
    = 2D\left(
        \frac{\partial m_x}{\partial x}
        + \frac{\partial m_y}{\partial y}
    \right)
\end{equation}
%
Applying~\eqref{eq:heff_def}:
%
\begin{equation}
\Hdmi = \frac{2D}{\Ms} \left(
    \frac{\partial m_z}{\partial x}\,\hat{x}
    + \frac{\partial m_z}{\partial y}\,\hat{y}
    - \left(
        \frac{\partial m_x}{\partial x}
        + \frac{\partial m_y}{\partial y}
    \right)\zhat
\right)
\end{equation}

%----------------------------------------------------------------------
\subsection{Derivation of the Anisotropy Field}

From the anisotropy energy $\mathcal{H}_{\text{anis}} = -K \int
m_z^2 \, dV$:
%
\begin{equation}
\delta\mathcal{H}_{\text{anis}} = -2K \int
    m_z \, \delta m_z \; dV
    = -2K \int (m_z \, \zhat) \cdot \delta\mb{m} \; dV
\end{equation}
%
Therefore:
%
\begin{equation}
\Hanis = \frac{2K}{\Ms} \, m_z \, \zhat
\end{equation}
%
The demagnetization field follows identically with
$K \to -\mu_0\Ms^2/2$, giving
$\Hdemag = -\mu_0\Ms \, m_z \, \zhat$, and the combined field uses
$\Keff = K - \mu_0\Ms^2/2$.

%----------------------------------------------------------------------
\subsection{Derivation of the Zeeman and RKKY Fields}

The Zeeman energy
$\mathcal{H}_{\text{Z}} = -\Ms \int \mb{m} \cdot
\mb{B}_{\text{ext}} \, dV$ is linear in $\mb{m}$:
%
\begin{equation}
\delta\mathcal{H}_{\text{Z}} = -\Ms \int
    \mb{B}_{\text{ext}} \cdot \delta\mb{m} \; dV
\qquad\Longrightarrow\qquad
\Hext = \mb{B}_{\text{ext}}
\end{equation}

Similarly, the RKKY energy
$\mathcal{H}_{\text{RKKY}} = J_{\text{RKKY}} \int
\mb{m}_{\text{top}} \cdot \mb{m}_{\text{bot}} \, dV$ gives,
when varying $\mb{m}_{\text{top}}$:
%
\begin{equation}
\left.\Hrkky\right|_{\text{top}} =
    -\frac{J_{\text{RKKY}}}{\Ms}\,\mb{m}_{\text{bot}}
    = -H_{\text{RKKY}}\,\mb{m}_{\text{bot}}
\end{equation}
%
and symmetrically for the bottom layer.

%----------------------------------------------------------------------
\subsection{Physical Interpretation of the LLG Equation}

With the effective field defined, the LLG equation
\begin{equation}
\frac{d\mb{m}}{dt} =
    \underbrace{-\gamma\,\mb{m} \times \Heff}_{\text{precession}}
    + \underbrace{\alpha\,\mb{m} \times
        \frac{d\mb{m}}{dt}}_{\text{damping}}
\end{equation}
has a clear energetic interpretation:

\begin{itemize}[nosep]
\item \textbf{Precession term}
    $(-\gamma\,\mb{m}\times\Heff)$:
    The torque is perpendicular to both $\mb{m}$ and $\Heff$,
    so the magnetization precesses around the effective field.
    Since $\mb{m} \cdot d\mb{m}/dt = 0$ (from $|\mb{m}|=1$)
    and the torque does no work on the field
    ($\Heff \cdot (\mb{m}\times\Heff) = 0$), this term
    \textbf{conserves energy}.

\item \textbf{Damping term}
    $(\alpha\,\mb{m}\times d\mb{m}/dt)$:
    This produces a torque component that tilts $\mb{m}$ toward
    $\Heff$.  One can show that the energy dissipation rate is:
    %
    \begin{equation}
    \frac{d\mathcal{H}}{dt} =
        -\frac{\alpha\Ms}{\gamma}
        \int \left|\frac{d\mb{m}}{dt}\right|^2 dV
        \leq 0
    \end{equation}
    %
    Energy is always dissipated ($d\mathcal{H}/dt \leq 0$),
    driving the system toward the energy minimum.
\end{itemize}

Together, the two terms describe a \textbf{damped precession}:
the magnetization spirals around the effective field while
gradually aligning with it.  The Hamiltonian defines the energy
landscape; the effective field is the local gradient of that
landscape; and the LLG equation is the equation of motion that
navigates the magnetization toward equilibrium.

The Slonczewski spin-transfer torque terms (Section~6) add a
non-conservative driving force from the electric current, which
can sustain steady-state motion (e.g.\ skyrmion translation)
against the damping.

%======================================================================
\section{Landau--Lifshitz--Gilbert--Slonczewski Equation (LLGS)}
%======================================================================

The magnetization dynamics of each layer are governed by the LLGS
equation. For the unit magnetization vector
$\mb{m} = \mb{M}/\Ms$:
%
\begin{equation}\label{eq:llgs}
\frac{d\mb{m}}{dt} =
    -\gamma \, \mb{m} \times \Heff
    + \alpha \, \mb{m} \times \frac{d\mb{m}}{dt}
    + \gamma H_{\text{DL}}
        \frac{\mb{m} \times (\phat \times \mb{m})}
             {|\phat \times \mb{m}|}
    + \gamma H_{\text{FL}}
        \frac{\phat \times \mb{m}}
             {|\phat \times \mb{m}|}
\end{equation}
%
where:
\begin{itemize}[nosep]
\item $\gamma$ is the gyromagnetic ratio (rad/s/T),
\item $\alpha$ is the Gilbert damping constant,
\item $\Heff$ is the effective magnetic field (Tesla),
\item $H_{\text{DL}} = \chi_{\text{DL}} \cdot J$ is the damping-like
    (DL) SOT field,
\item $H_{\text{FL}} = \chi_{\text{FL}} \cdot J$ is the field-like
    (FL) SOT field,
\item $\phat$ is the spin polarization direction
    ($\phat = \zhat \times \hat{J}$),
\item $\mb{J} = J\,\hat{J}$ is the current density: $J$ is the
    magnitude (A/m$^2$), $\hat{J}$ is the unit vector along the
    current flow,
\item $\zhat$ is the film-normal unit vector (out-of-plane).
\end{itemize}

%----------------------------------------------------------------------
\subsection{Explicit Form}
\label{sec:llgs_explicit_derivation}

The implicit form~\eqref{eq:llgs} contains $d\mb{m}/dt$ on both sides
through the Gilbert damping term. To convert it to a fully explicit
ODE solvable by Runge--Kutta, we eliminate the right-hand-side
$d\mb{m}/dt$ algebraically. The standard procedure
follows Mallinson~\cite{mallinson1987} for the LLG part and
Khvalkovskiy~\emph{et al.}~\cite{khvalkovskiy2013} for the SOT
extension.

\paragraph{Step 1: Compact form.}
Define $\mb{s} = (\phat\times\mb{m})/|\phat\times\mb{m}|$ and group
the right-hand side of~\eqref{eq:llgs} as
%
\begin{equation}\label{eq:F_def}
\frac{d\mb{m}}{dt} = \mb{F} + \alpha\,\mb{m}\times\frac{d\mb{m}}{dt},
\qquad
\mb{F} = -\gamma\,\mb{m}\times\Heff
        + \gamma H_{\text{DL}}\,(\mb{m}\times\mb{s})
        + \gamma H_{\text{FL}}\,\mb{s}.
\end{equation}

\paragraph{Step 2: Cross with $\mb{m}$.}
Take $\mb{m}\times$ of both sides of~\eqref{eq:F_def}:
%
\begin{equation}\label{eq:cross_m}
\mb{m}\times\frac{d\mb{m}}{dt}
    = \mb{m}\times\mb{F}
        + \alpha\,\mb{m}\times\Big(\mb{m}\times\frac{d\mb{m}}{dt}\Big).
\end{equation}
%
The triple product is simplified by the BAC--CAB identity
$\mb{a}\times(\mb{b}\times\mb{c}) = \mb{b}(\mb{a}\cdot\mb{c})
- \mb{c}(\mb{a}\cdot\mb{b})$. Using $|\mb{m}|=1$ (so
$\mb{m}\cdot\mb{m}=1$) and
$\mb{m}\cdot d\mb{m}/dt = \tfrac{1}{2}d|\mb{m}|^2/dt = 0$:
%
\begin{equation}\label{eq:bac_cab_mdot}
\mb{m}\times\Big(\mb{m}\times\frac{d\mb{m}}{dt}\Big)
    = \mb{m}\Big(\mb{m}\cdot\frac{d\mb{m}}{dt}\Big)
    - \frac{d\mb{m}}{dt}(\mb{m}\cdot\mb{m})
    = -\frac{d\mb{m}}{dt}.
\end{equation}
%
Substituting~\eqref{eq:bac_cab_mdot} into~\eqref{eq:cross_m}:
%
\begin{equation}\label{eq:cross_m_solved}
\mb{m}\times\frac{d\mb{m}}{dt}
    = \mb{m}\times\mb{F} - \alpha\,\frac{d\mb{m}}{dt}.
\end{equation}

\paragraph{Step 3: Solve for $d\mb{m}/dt$.}
Insert~\eqref{eq:cross_m_solved} back into the
$\alpha\,\mb{m}\times d\mb{m}/dt$ term of~\eqref{eq:F_def}:
%
\begin{equation}
\frac{d\mb{m}}{dt}
    = \mb{F} + \alpha\big(\mb{m}\times\mb{F} - \alpha\,\tfrac{d\mb{m}}{dt}\big),
\end{equation}
%
collect the $d\mb{m}/dt$ terms on the left, and divide by
$1+\alpha^2$:
%
\begin{equation}\label{eq:llgs_F_explicit}
\boxed{\;
\frac{d\mb{m}}{dt}
    = \frac{1}{1+\alpha^2}\Big[\mb{F} + \alpha\,\mb{m}\times\mb{F}\Big].
\;}
\end{equation}

\paragraph{Step 4: Expand $\mb{m}\times\mb{F}$.}
The three contributions to $\mb{F}$ in~\eqref{eq:F_def} give:
%
\begin{align}
\mb{m}\times\big(-\gamma\mb{m}\times\Heff\big)
    &= -\gamma\,\mb{m}\times(\mb{m}\times\Heff), \\
\mb{m}\times\big(\gamma H_{\text{DL}}\,\mb{m}\times\mb{s}\big)
    &= \gamma H_{\text{DL}}\,\mb{m}\times(\mb{m}\times\mb{s})
     = -\gamma H_{\text{DL}}\,\mb{s},
        \label{eq:DL_simplify}\\
\mb{m}\times\big(\gamma H_{\text{FL}}\,\mb{s}\big)
    &= \gamma H_{\text{FL}}\,\mb{m}\times\mb{s},
\end{align}
%
where~\eqref{eq:DL_simplify} again uses BAC--CAB with
$\mb{m}\cdot\mb{s}=0$ (since $\mb{s}\propto\phat\times\mb{m}$ is
perpendicular to $\mb{m}$):
$\mb{m}\times(\mb{m}\times\mb{s}) = \mb{m}(\mb{m}\cdot\mb{s})
- \mb{s} = -\mb{s}$.

\paragraph{Step 5: Collect coefficients.}
Assembling $\mb{F} + \alpha\mb{m}\times\mb{F}$ and dividing by
$1+\alpha^2$:
%
\begin{equation}\label{eq:llgs_explicit}
\boxed{\;
\frac{d\mb{m}}{dt}
    = \frac{\gamma}{1 + \alpha^2} \Big[
        -\mb{m}\times\Heff
        - \alpha\,\mb{m}\times(\mb{m}\times\Heff)
        + (H_{\text{DL}} + \alpha H_{\text{FL}})\,(\mb{m}\times\mb{s})
        + (H_{\text{FL}} - \alpha H_{\text{DL}})\,\mb{s}
    \Big].
\;}
\end{equation}

The coefficients $H_{\text{DL}} + \alpha H_{\text{FL}}$ and
$H_{\text{FL}} - \alpha H_{\text{DL}}$ are the standard
``mixed'' SOT coefficients of the explicit
form~\cite{khvalkovskiy2013}; the cross-coupling between DL and FL
through $\alpha$ arises entirely from
Step~3's $\alpha\,\mb{m}\times\mb{F}$ term and vanishes at $\alpha=0$.
%
The implementation in \texttt{src/simulator/integrator.py}
(\texttt{llgs\_rhs}) precomputes
$\gamma' = \gamma/(1+\alpha^2)$,
$c_{\text{DL}} = H_{\text{DL}} + \alpha H_{\text{FL}}$,
$c_{\text{FL}} = H_{\text{FL}} - \alpha H_{\text{DL}}$,
and evaluates~\eqref{eq:llgs_explicit} directly without further
algebra.

%======================================================================
\section{Effective Field Contributions}
%======================================================================

The total effective field for each layer is:
%
\begin{equation}
\Heff = \Hex + \Hdmi + \Hanis + \Hext + \Hrkky + \Hdemag
\end{equation}
%
All fields are computed in Tesla.

%----------------------------------------------------------------------
\subsection{Exchange Field}

On a 2D square lattice with lattice constant $a$, the discrete
exchange field is the lattice Laplacian:
%
\begin{equation}
\Hex = \frac{2\Aex}{\Ms \, a^2}
    \left(
        \mb{m}_{+x} + \mb{m}_{-x}
        + \mb{m}_{+y} + \mb{m}_{-y}
        - 4\mb{m}
    \right)
\end{equation}

%----------------------------------------------------------------------
\subsection{Dzyaloshinskii--Moriya Interaction}

The discrete interfacial (N\'{e}el-type) DMI field:
%
\begin{align}
H_{\text{DMI},x} &= \frac{D}{\Ms \, a}
    (m_{+x,z} - m_{-x,z}) \\
H_{\text{DMI},y} &= \frac{D}{\Ms \, a}
    (m_{+y,z} - m_{-y,z}) \\
H_{\text{DMI},z} &= -\frac{D}{\Ms \, a}
    (m_{+x,x} - m_{-x,x} + m_{+y,y} - m_{-y,y})
\end{align}
%
where $D > 0$ favors left-handed N\'{e}el skyrmions, consistent with
the Pt/Co interface.

%----------------------------------------------------------------------
\subsection{Perpendicular Magnetic Anisotropy}

Uniaxial anisotropy along the film normal ($\zhat$):
%
\begin{equation}
\Hanis = \frac{2K}{\Ms} \, m_z \, \zhat
\end{equation}

%----------------------------------------------------------------------
\subsection{Thin-Film Demagnetization}

In a thin film, the demagnetizing field is:
%
\begin{equation}
\Hdemag = -\mu_0 \Ms \, m_z \, \zhat
\end{equation}

The combined anisotropy + demagnetization field uses the effective
anisotropy:
%
\begin{equation}
\Keff = K - \frac{\mu_0 \Ms^2}{2}
\qquad\Longrightarrow\qquad
\mb{H}_{\text{anis+demag}} =
    \frac{2\Keff}{\Ms} \, m_z \, \zhat
\end{equation}

The system in~\cite{pham2024} is close to the spin reorientation
transition ($\Keff \approx 0$), which favors skyrmion formation.

%----------------------------------------------------------------------
\subsection{Zeeman Field}

A uniform external magnetic field $\Hext$ (in Tesla), identical at all
lattice sites.

%----------------------------------------------------------------------
\subsection{RKKY Interlayer Coupling}

The antiferromagnetic RKKY interaction promotes antiparallel alignment:
%
\begin{align}
\mb{H}_{\text{RKKY, on top}} &=
    -H_{\text{RKKY}} \, \mb{m}_{\text{bot}} \\
\mb{H}_{\text{RKKY, on bot}} &=
    -H_{\text{RKKY}} \, \mb{m}_{\text{top}}
\end{align}

%======================================================================
\section{Spin-Orbit Torques}
%======================================================================

The spin-orbit torques (SOT) are the mechanism by which an electric
current drives skyrmion motion.  They originate from the
\textbf{spin Hall effect} (SHE) in the heavy-metal Pt layers
adjacent to the ferromagnetic Co.

\subsection{Spin Hall Effect}

When an in-plane charge current density $\mb{J}$ flows through the
Pt layer, the strong spin-orbit coupling deflects spin-up and
spin-down electrons in opposite transverse directions.  This
generates a transverse \emph{spin current} that accumulates at the
Pt/Co interface.  The spin polarization direction is perpendicular
to both the current and the film normal:
%
\begin{equation}
\phat = \zhat \times \hat{J}
\end{equation}
%
For current along $\hat{x}$, the spin polarization is
$\phat = \hat{y}$.  The efficiency of the conversion from charge
current to spin current is characterized by the spin Hall angle
$\theta_{\text{SH}}$, which is encapsulated in the experimentally
measured SOT coefficients $\chi_{\text{DL}}$ and $\chi_{\text{FL}}$.

\subsection{SOT Torques in the LLGS Equation}

The spin current accumulated at the interface exerts two distinct
torques on the Co magnetization, which enter as the Slonczewski
terms in the LLGS equation:
%
\begin{itemize}[nosep]
\item \textbf{Damping-like (DL) torque}:
    $\tau_{\text{DL}} = H_{\text{DL}} \,
    \mb{m} \times (\phat \times \mb{m})
    / |\phat \times \mb{m}|$.
    This acts like an additional damping (or anti-damping) that
    pushes the magnetization toward or away from $\phat$.  It is
    the dominant torque driving skyrmion motion.

\item \textbf{Field-like (FL) torque}:
    $\tau_{\text{FL}} = H_{\text{FL}} \,
    (\phat \times \mb{m})
    / |\phat \times \mb{m}|$.
    This is equivalent to an effective magnetic field along $\phat$
    and is typically smaller in magnitude.
\end{itemize}

The SOT effective field magnitudes are proportional to the applied
current density:
%
\begin{equation}
H_{\text{DL}} = \chi_{\text{DL}} \cdot J,
\qquad
H_{\text{FL}} = \chi_{\text{FL}} \cdot J
\end{equation}
%
where $\chi_{\text{DL}}$ and $\chi_{\text{FL}}$ are the SOT
coefficients (in T\,A$^{-1}$\,m$^2$) measured experimentally.
These coefficients encapsulate the spin Hall angle, the spin
diffusion length, and the interface spin transparency into a single
effective parameter.

\subsection{SOT Terms in the Explicit LLGS}

In the explicit form of the LLGS equation~\eqref{eq:llgs_explicit},
the SOT terms appear as the last two contributions:
%
\begin{equation}
\frac{d\mb{m}}{dt} = \frac{\gamma}{1+\alpha^2} \Big[
    \underbrace{
        -\mb{m} \times \mb{H}
        - \alpha\,\mb{m} \times (\mb{m} \times \mb{H})
    }_{\text{precession + damping}}
    + \underbrace{
        (H_{\text{DL}} + \alpha H_{\text{FL}})
            (\mb{m} \times \mb{s})
        + (H_{\text{FL}} - \alpha H_{\text{DL}})\,\mb{s}
    }_{\text{SHE-driven SOT}}
\Big]
\end{equation}
%
where $\mb{s} = (\phat \times \mb{m})/|\phat \times \mb{m}|$.
Without the SHE ($J = 0$), the SOT terms vanish and the skyrmion
relaxes to its equilibrium configuration.  With nonzero current,
the DL-SOT exerts a net force on the skyrmion domain wall, driving
translational motion.  In the SAF geometry, the opposite topological
charges of the two layers cancel the Magnus force, resulting in
motion along the current direction with no skyrmion Hall
effect~\cite{pham2024}.

%======================================================================
\section{Magnus Force and the Thiele Equation}
%======================================================================

The Magnus force is not a separate term in the LLGS equation --- it
\emph{emerges} from it when skyrmion motion is analyzed as a
collective coordinate problem.  If the skyrmion moves rigidly with
velocity $\mb{v}$, integrating the LLGS torque equation over the
skyrmion area yields the Thiele equation~(42):
%
\begin{equation}\label{eq:thiele}
\mb{F}_{\text{SOT}} + \mb{G} \times \mb{v}
    + \alpha [\mb{D}] \mb{v} = 0
\end{equation}
%
The three terms represent:
%
\begin{itemize}[nosep]
\item $\mb{F}_{\text{SOT}}$: the driving force from the spin-orbit
    torque, directed along the current.
\item $\mb{G} \times \mb{v}$: the \textbf{Magnus force}, directed
    perpendicular to the skyrmion velocity.  The gyrovector is
    $\mb{G} = G\,\zhat$ with $G = 4\pi Q$, where $Q$ is the
    topological charge.  This force is analogous to the Magnus
    effect on a spinning ball and is a direct consequence of the
    skyrmion's non-trivial topology.
\item $\alpha[\mb{D}]\mb{v}$: the dissipative drag from Gilbert
    damping, where $[\mb{D}]$ is the dissipation tensor.
\end{itemize}
%
The Magnus force deflects the skyrmion sideways relative to the
driving force --- this is the \textbf{skyrmion Hall effect}.  From
\eqref{eq:thiele}, the skyrmion Hall angle is:
%
\begin{equation}
\tan\theta_{\text{SkHE}} = \frac{v_y}{v_x}
    = \frac{G}{\alpha D}
\end{equation}
%
where $D$ is the diagonal component of the dissipation tensor.

In a \textbf{single ferromagnetic layer} with $Q = \pm 1$, the
gyrovector $G = \pm 4\pi \neq 0$, and the skyrmion is deflected
at a significant angle to the current (experimentally up to
$60^\circ$--$85^\circ$ at high velocities~\cite{pham2024}).

In the \textbf{SAF}, the two layers carry opposite topological
charges ($Q_{\text{top}} = -Q_{\text{bot}}$), so their gyrovectors
cancel:
%
\begin{equation}
G_{\text{net}} = G_{\text{top}} + G_{\text{bot}}
    = 4\pi Q_{\text{top}} + 4\pi Q_{\text{bot}} = 0
\end{equation}
%
The Magnus forces on the two layers point in opposite directions
and cancel exactly.  The skyrmion moves straight along the current
direction with $\theta_{\text{SkHE}} = 0$.  This is the key
advantage of the SAF geometry for skyrmion-based devices, as
confirmed both experimentally
($\theta_{\text{SkHE}} = 6.4^\circ \pm 16.4^\circ$) and in our
simulations ($\theta_{\text{SkHE}} = 0.3^\circ$).

%======================================================================
\section{Topological Charge}
%======================================================================

The topological charge (skyrmion number) is:
%
\begin{equation}
Q = \frac{1}{4\pi} \int \mb{m} \cdot \left(
    \frac{\partial \mb{m}}{\partial x} \times
    \frac{\partial \mb{m}}{\partial y}
\right) dx \, dy
\end{equation}
%
For a single skyrmion, $Q = \pm 1$. In the SAF, the two layers carry
opposite topological charges
($Q_{\text{top}} = -Q_{\text{bot}}$), leading to a net topological
charge of zero and hence no skyrmion Hall effect.

%======================================================================
\section{Skyrmion Profile (Initial Condition)}
%======================================================================

The initial N\'{e}el skyrmion profile is parametrized by:
%
\begin{equation}
\theta(r) = \begin{cases}
    \pi \left(1 - \dfrac{r}{R}\right) & r < R \\[6pt]
    0 & r \geq R
\end{cases}
\end{equation}
%
\begin{equation}
m_x = \sin\theta \, \cos\varphi, \quad
m_y = \sin\theta \, \sin\varphi, \quad
m_z = \cos\theta
\end{equation}
%
where $r$ is the distance from the skyrmion center, $R$ is the
skyrmion radius, and
$\varphi = \text{atan2}(y - y_0, x - x_0)$ is the azimuthal angle
(N\'{e}el helicity).

The top layer has core down ($m_z = -1$ at center), the bottom layer
has core up ($m_z = +1$ at center), consistent with antiferromagnetic
RKKY coupling.

%======================================================================
\section{Time Integration}
%======================================================================

The explicit LLGS equation~\eqref{eq:llgs_explicit} is integrated
using the classical 4th-order Runge--Kutta (RK4) method. At each
substep, the spin configuration is renormalized to enforce
$|\mb{m}| = 1$:
%
\begin{align}
\mb{k}_1 &= \Delta t \, f(\mb{m}) \\
\mb{k}_2 &= \Delta t \, f\!\left(
    \text{norm}\!\left(
        \mb{m} + \tfrac{\mb{k}_1}{2}
    \right)\right) \\
\mb{k}_3 &= \Delta t \, f\!\left(
    \text{norm}\!\left(
        \mb{m} + \tfrac{\mb{k}_2}{2}
    \right)\right) \\
\mb{k}_4 &= \Delta t \, f\!\left(
    \text{norm}\!\left(
        \mb{m} + \mb{k}_3
    \right)\right) \\
\mb{m}^{n+1} &= \text{norm}\!\left(
    \mb{m}^n + \frac{
        \mb{k}_1 + 2\mb{k}_2 + 2\mb{k}_3 + \mb{k}_4
    }{6}\right)
\end{align}
%
Both SAF layers are evolved simultaneously, as the RKKY coupling
makes them interdependent.

%======================================================================
\section{Boundary Conditions}
%======================================================================

Periodic boundary conditions (PBC) are applied in both $x$ and $y$
directions, implemented via \texttt{np.roll} on the spin arrays.

%======================================================================
\section{Simulation Parameters}
%======================================================================

All material parameters are taken from~\cite{pham2024} for the
optimized SAF stack
Pt(3)/Co(1.58)/Ru(0.85)/Pt(0.5)/Co(1.58)/Ru(0.85)
(thicknesses in nm).

\begin{table}[h!]
\centering
\caption{Material parameters from~\cite{pham2024}.}
\begin{tabular}{@{}llll@{}}
\toprule
Parameter & Symbol & Value & Unit \\
\midrule
Saturation magnetization
    & $\Ms$ & $1.43 \times 10^6$ & A/m \\
Exchange stiffness
    & $\Aex$ & $16 \times 10^{-12}$ & J/m \\
DMI constant
    & $D$ & $0.62 \times 10^{-3}$ & J/m$^2$ \\
Anisotropy (top layer)
    & $K_{\text{top}}$ & $1.294 \times 10^6$ & J/m$^3$ \\
Anisotropy (bottom layer)
    & $K_{\text{bot}}$ & $1.31 \times 10^6$ & J/m$^3$ \\
Gilbert damping
    & $\alpha$ & 0.14 & -- \\
Co layer thickness
    & $\tCo$ & 1.3 & nm \\
Gyromagnetic ratio
    & $\gamma$ & $194.8 \times 10^9$ & rad/(s\,T) \\
RKKY coupling field
    & $H_{\text{RKKY}}$ & 205 & mT \\
DL-SOT coefficient
    & $\chi_{\text{DL}}$
    & $2.21 \times 10^{-14}$
    & T\,A$^{-1}$\,m$^2$ \\
FL-SOT coefficient
    & $\chi_{\text{FL}}$
    & $0.53 \times 10^{-14}$
    & T\,A$^{-1}$\,m$^2$ \\
Domain wall width
    & $\Delta$ & 27 & nm \\
\bottomrule
\end{tabular}
\end{table}

\subsection*{Derived quantities}

\begin{table}[h!]
\centering
\begin{tabular}{@{}lll@{}}
\toprule
Quantity & Formula & Value \\
\midrule
$K_{\text{eff, top}}$
    & $K_{\text{top}} - \mu_0 \Ms^2 / 2$
    & $\approx 9.2 \times 10^3$ J/m$^3$ \\
$K_{\text{eff, bot}}$
    & $K_{\text{bot}} - \mu_0 \Ms^2 / 2$
    & $\approx 2.5 \times 10^4$ J/m$^3$ \\
Domain wall width
    & $\Delta = \sqrt{\Aex / \Keff}$
    & $\approx 30.5$ nm \\
Skyrmion diameter
    & (from paper formula)
    & $\approx 197$ nm (paper), 160 nm (sim) \\
\bottomrule
\end{tabular}
\end{table}

\subsection*{Numerical parameters}

\begin{table}[h!]
\centering
\begin{tabular}{@{}lll@{}}
\toprule
Parameter & Value & Description \\
\midrule
$n_x \times n_y$ & $256 \times 256$ & Lattice sites \\
$a$ & 2 nm & Lattice constant \\
$\Delta t$ & $5 \times 10^{-14}$ s & Time step \\
Relaxation steps & 10\,000 & 500 ps at $J = 0$ \\
Drive steps & 5\,000 & 250 ps with current \\
Current density & $4 \times 10^{11}$ A/m$^2$ & \\
Spin polarization & $\phat = \hat{y}$
    & From spin Hall effect \\
Initial skyrmion radius & 80 nm
    & Relaxes to equilibrium \\
Domain wall width (init.) & 27 nm
    & From paper (Table S2) \\
\bottomrule
\end{tabular}
\end{table}

\subsection*{Stability condition}

The RK4 time step must satisfy:
%
\begin{equation}
\omega_{\text{max}} \cdot \Delta t < 2.8
\end{equation}
%
where $\omega_{\text{max}} = \gamma \cdot H_{\text{max}}$ and
$H_{\text{max}} \approx 4 C_{\text{ex}} \approx 90$~T (exchange
field at the skyrmion core boundary). With the chosen parameters,
$\omega_{\text{max}} \cdot \Delta t \approx 0.87$, within the RK4
stability region.

%======================================================================
\section{Simulation Results vs.\ Paper}
%======================================================================

\begin{table}[h!]
\centering
\caption{Comparison of simulation results with~\cite{pham2024}.}
\begin{tabular}{@{}lll@{}}
\toprule
Quantity & Simulation & Paper \\
\midrule
Equilibrium diameter & 160 nm & 197 nm \\
Velocity ($J = 4 \times 10^{11}$ A/m$^2$)
    & 500 m/s & $\sim$400 m/s (micromagnetic) \\
Skyrmion Hall angle & $0.3^\circ$ & $\sim 0^\circ$ (SAF) \\
Topological charge & $\pm 0.999$ & $\pm 1$ \\
\bottomrule
\end{tabular}
\end{table}

The diameter discrepancy (160 vs 197~nm) arises from the thin-film
demagnetization approximation (local $N_z = 1$) compared to the full
magnetostatic FFT calculation in MuMax3, and from parameter
uncertainties ($D$ has $\pm 39\%$ error bars, $K$ has $\pm 7\%$).
The velocity is within the uncertainty band of the damping constant
($\alpha = 0.14 \pm 0.04$). The four approximations are detailed in
Sec.~\ref{sec:implementation_notes}.

%======================================================================
\section{Implementation Notes}
\label{sec:implementation_notes}
%======================================================================

Four sources of model--experiment disagreement enter the
simulation--paper comparison of Sec.~13. They are intentional
modeling choices or material-parameter uncertainties, not bugs.

%----------------------------------------------------------------------
\subsection{Local $N_z = 1$ demag vs full FFT demag}

The micromagnetic demagnetization tensor $\hat N(\mb{k})$ depends on
wavevector. For a thin slab of thickness $\tCo$:
%
\begin{equation}
N_{zz}(k) = f(|k|\,\tCo),
\qquad f(x) = \frac{1 - e^{-x}}{x}.
\end{equation}
%
\begin{itemize}[leftmargin=2em, nosep]
    \item $|k|\,\tCo \to 0$ (long wavelength, uniform film):
        $f \to 1$, $N_{zz} \to 1$, $N_{xx, yy} \to 0$. In real space:
        $\Hdemag = -\mu_0 \Ms\, m_z\,\zhat$ --- a \textbf{local}
        field driven only by $m_z$ at the same point.
    \item $|k|\,\tCo \to \infty$ (short wavelength): $f \to 0$, no
        demag.
\end{itemize}

The \textbf{local $N_z = 1$ approximation} assumes $N_{zz}=1$
everywhere and absorbs the demag into the anisotropy:
%
\begin{equation}
\Keff = K - \frac{1}{2}\mu_0 \Ms^2.
\end{equation}
%
The main RK4 loop in \texttt{src/simulator/main.py} uses this
approximation; it never calls \texttt{demag.py}. It is dimensionally
correct only when:
\begin{enumerate}[leftmargin=2em, nosep]
    \item the texture varies slowly compared to $\tCo$ --- true here
        ($R \sim 80$~nm vs $\tCo = 1.3$~nm, $|k|\,\tCo \sim 0.016$);
    \item the magnetization is dominantly out-of-plane.
\end{enumerate}
The second assumption \textbf{fails inside the domain wall}, where
$\mb{m}$ is in-plane, creates surface magnetic charges, and induces
non-local stray fields between in-plane spins.

The full FFT demag (in \texttt{src/simulator/demag.py}) computes
$\hat N(\mb{k})$ exactly via FFT, including in-plane components and
the inter-layer dipolar kernel across the Ru spacer. MuMax3 does
this by default; the simulator's main loop does not. This is the
\textbf{dominant source of the 160 nm (sim) vs 197 nm (paper)
diameter mismatch}: non-local demag lowers the in-plane DW energy,
allowing a wider wall and a larger skyrmion.

%----------------------------------------------------------------------
\subsection{DMI uncertainty: $D \pm 39\%$}

Interfacial DMI cannot be measured directly. Standard techniques:
\begin{itemize}[leftmargin=2em, nosep]
    \item \textbf{Brillouin light scattering (BLS):} asymmetric
        spin-wave dispersion $\omega(+k)\neq\omega(-k)$ caused by DMI.
        Sensitive to interface roughness and Pt thickness.
    \item \textbf{Domain-wall creep velocity vs in-plane field:}
        $v(H_x)$ asymmetry yields $D$.
    \item \textbf{Skyrmion size vs out-of-plane field:} equilibrium
        $R$ is monotone in $D/D_c$.
\end{itemize}
Each method has 10--30\% systematic uncertainty (interface quality,
sample variation, fitting models). Cross-method agreement is often
factor-of-2. Pham~\emph{et al.}~\cite{pham2024} quote
$D = 0.62 \pm 0.24$~mJ/m$^2$ ($\approx \pm 39\%$).

The skyrmion size depends strongly on $D$: near the Bogdanov--Hubert
threshold $D_c = (4/\pi)\sqrt{\Aex\,\Keff}$, the equilibrium $R$
diverges. A $\pm 39\%$ variation in $D$ shifts $D/D_c$ enough to
change $R$ by tens of nm.

%----------------------------------------------------------------------
\subsection{Anisotropy uncertainty: $K \pm 7\%$}

$K$ is measured via FMR resonance fields, anomalous-Hall hysteresis,
or MOKE saturation fields. These methods are well-standardized and
agree within 5--10\% between labs. Hence $\pm 7\%$ on $K$.

The catch is that $\Keff = K - \mu_0 \Ms^2/2$ is a
\textbf{small difference of large numbers}:
%
\begin{equation}
K = 1.294\times 10^6,\quad
\mu_0 \Ms^2/2 = 1.286\times 10^6\ \text{J/m}^3
\;\Longrightarrow\;
\Keff \approx 8\times 10^3\ \text{J/m}^3.
\end{equation}
%
A 7\% error on $K$ ($\sim 9\times 10^4$~J/m$^3$) is
\textbf{10$\times$ larger than $\Keff$ itself} --- the system sits in
the spin reorientation transition regime where $\Keff$ is barely
positive. The domain-wall width $\Delta = \sqrt{\Aex/\Keff}$
depends on $1/\sqrt{\Keff}$, so few-percent errors on $K$ propagate
to large errors on $\Delta$ and the skyrmion size.

%----------------------------------------------------------------------
\subsection{Gilbert damping uncertainty: $\alpha \pm 0.04$}

Damping in Co/Pt is enhanced by \textbf{spin pumping}: precessing
magnetization in Co injects a transverse spin current into Pt, where
strong spin-orbit scattering dissipates it. This adds an interfacial
term to the bulk Co damping. Measurements:
\begin{itemize}[leftmargin=2em, nosep]
    \item \textbf{FMR linewidth} $\Delta H$ vs frequency $f$:
        slope $= \alpha \cdot 2\pi/\gamma$.
    \item \textbf{Time-resolved MOKE} precession decay.
\end{itemize}
Both give $\alpha \approx 0.14 \pm 0.04$ for the Co/Pt stack
($\sim \pm 30\%$).

In the SAF, the Magnus term cancels between layers
(Sec.~7), so the Thiele equation
reduces to a damping-limited drift:
%
\begin{equation}
v \propto \frac{1}{\alpha}.
\end{equation}
%
A $\pm 0.04$ uncertainty in $\alpha$ around 0.14 propagates to
$\pm 30\%$ in velocity. The simulation gives 500~m/s, the MuMax3
reference gives $\sim$400~m/s --- a 25\% gap, comfortably inside the
damping uncertainty band.

%----------------------------------------------------------------------
\subsection{Summary}

\begin{table}[h!]
\centering
\begin{tabular}{@{}llll@{}}
\toprule
Approximation & Mechanism & Affects & Magnitude \\
\midrule
Local $N_z = 1$ demag
    & non-local in-plane dipolar ignored
    & DW energy, skyrmion size
    & $\sim$20\% on diameter \\
$D \pm 39\%$
    & DMI hard to measure
    & skyrmion size, $D_c$
    & exponential near threshold \\
$K \pm 7\%$
    & material spread
    & $\Keff$ via subtraction
    & amplified $\sim$10$\times$ on $\Keff$ \\
$\alpha \pm 0.04$
    & spin pumping enhancement
    & drift velocity $\propto 1/\alpha$
    & $\sim$30\% on velocity \\
\bottomrule
\end{tabular}
\end{table}

%======================================================================
\section{Output Format}
%======================================================================

Spin configurations are written in LAMMPS dump format. Since OVITO
does not natively support spin degrees of freedom, the three spin
components ($m_x$, $m_y$, $m_z$) are written as force components
(\texttt{fx}, \texttt{fy}, \texttt{fz}), which OVITO can render as
vector arrows. Positions are output in nanometers.

Type~1 corresponds to the top Co layer ($z = 0$), type~2 to the
bottom Co layer ($z = -\tCo$).

%======================================================================
\appendix
\section{Derivation of the Domain Wall Profile}
\label{app:dw}
%======================================================================

The skyrmion initial condition uses the one-dimensional domain wall
profile obtained from the Euler--Lagrange equation.  This appendix
derives that profile from first principles.

\subsection*{From 3D energy to 1D energy density}

The starting point is the three-dimensional micromagnetic energy for
a thin film with exchange and effective anisotropy (Sections~2.1
and~2.3--2.4):
%
\begin{equation}
E = \int_V \left[
    \Aex \left|\nabla \mb{m}\right|^2
    - \Keff \, m_z^2
\right] dV
\end{equation}
%
For a domain wall that varies only along~$x$ in a film of
thickness~$\tCo$ and cross-sectional area~$L_y \, \tCo$, the volume
integral reduces to:
%
\begin{equation}
E = L_y \, \tCo \int_{-\infty}^{+\infty} \left[
    \Aex \left|\frac{d\mb{m}}{dx}\right|^2
    - \Keff \, m_z^2
\right] dx
\end{equation}
%
We parametrize the unit magnetization by the polar angle
$\theta(x)$ and azimuthal angle $\phi(x)$:
%
\begin{equation}
\mb{m} = \big(
    \sin\theta\cos\phi, \;
    \sin\theta\sin\phi, \;
    \cos\theta
\big)
\end{equation}
%
Computing the gradient magnitude:
%
\begin{equation}
\left|\frac{d\mb{m}}{dx}\right|^2
    = \left(\frac{d\theta}{dx}\right)^{\!2}
    + \sin^2\!\theta
      \left(\frac{d\phi}{dx}\right)^{\!2}
\end{equation}
%
For a N\'{e}el wall (or Bloch wall), the azimuthal angle is
constant ($d\phi/dx = 0$), so the second term vanishes.
Substituting $m_z = \cos\theta$ and noting that
$-\Keff \cos^2\!\theta = -\Keff + \Keff \sin^2\!\theta$
(dropping the constant), the energy per unit cross-section becomes:
%
\begin{equation}
\frac{E}{L_y \, \tCo}
    = \int_{-\infty}^{+\infty} \varepsilon(x) \, dx
\end{equation}
%
with the one-dimensional energy density:
%
\begin{equation}\label{eq:dw_energy}
\boxed{\;
\varepsilon = \Aex
    \left(\frac{d\theta}{dx}\right)^{\!2}
    + \Keff \sin^2\!\theta
\;}
\end{equation}
%
The first term is the exchange cost of spatial variations in
$\theta$; the second is the anisotropy cost for deviating from the
easy-axis directions $\theta = 0$ ($m_z = +1$) or $\theta = \pi$
($m_z = -1$).

\subsection*{Euler--Lagrange equation}

Minimizing the total energy
$E = \int_{-\infty}^{+\infty} \varepsilon \, dx$
with respect to $\theta(x)$ gives the Euler--Lagrange equation:
%
\begin{equation}\label{eq:dw_EL}
2\Aex \frac{d^2\theta}{dx^2} = \Keff \sin(2\theta)
\end{equation}

\subsection*{First integral}

Multiplying both sides by $d\theta/dx$ and integrating once yields a
first integral (conservation of ``energy'' in the mechanical analogy):
%
\begin{equation}
\Aex \left(\frac{d\theta}{dx}\right)^{\!2}
    = \Keff \sin^2\!\theta
\end{equation}
%
which gives:
%
\begin{equation}\label{eq:dw_ode}
\frac{d\theta}{dx} = -\frac{\sin\theta}{\Delta},
\qquad
\Delta \equiv \sqrt{\frac{\Aex}{\Keff}}
\end{equation}
%
The negative sign ensures $\theta$ decreases from $\pi$ to $0$ as
$x$ increases (wall centered at $x = 0$).

\subsection*{Exact solution}

Separating variables in~\eqref{eq:dw_ode}:
%
\begin{equation}
\int \frac{d\theta}{\sin\theta} = -\frac{x}{\Delta} + C
\end{equation}
%
Using the identity
$\int \csc\theta \, d\theta
    = \ln|\tan(\theta/2)|$:
%
\begin{equation}
\ln\!\left|\tan\frac{\theta}{2}\right|
    = -\frac{x}{\Delta}
\end{equation}
%
where the integration constant is fixed by requiring $\theta(0)
= \pi/2$ (wall center). Solving for~$\theta$:
%
\begin{equation}\label{eq:dw_solution}
\boxed{\;
\theta(x) = 2\arctan\!\left(
    \exp\!\left(-\frac{x}{\Delta}\right)
\right)
\;}
\end{equation}

\subsection*{Properties}

The solution~\eqref{eq:dw_solution} has the following limits:
%
\begin{itemize}[nosep]
\item $x \to -\infty$: $\theta \to \pi$
    \quad ($m_z = -1$, down domain)
\item $x = 0$: $\theta = \pi/2$
    \quad ($m_z = 0$, wall center)
\item $x \to +\infty$: $\theta \to 0$
    \quad ($m_z = +1$, up domain)
\end{itemize}
%
The magnetization component $m_z = \cos\theta$ transitions
smoothly as:
%
\begin{equation}
m_z(x) = -\tanh\!\left(\frac{x}{\Delta}\right)
\end{equation}
%
confirming that $\Delta$ is the characteristic wall width.  Using
the parameters from~\cite{pham2024},
$\Delta = \sqrt{\Aex / \Keff} \approx 27$~nm.

\subsection*{Extension to a skyrmion}

For a radially symmetric skyrmion centered at the origin, the full
2D variational problem has no closed-form solution because the DMI
and the curvature of the domain wall add additional terms.  However,
the radial profile $\theta(r)$ is well approximated by wrapping the
1D wall into a circle at the skyrmion radius~$R_{sk}$:
%
\begin{equation}\label{eq:sk_profile}
\theta(r) = 2\arctan\!\left(
    \exp\!\left(-\frac{r - R_{sk}}{\Delta}\right)
\right)
\end{equation}
%
This gives:
%
\begin{itemize}[nosep]
\item $r \ll R_{sk}$: $\theta \to \pi$
    \quad (core, $m_z = -1$)
\item $r = R_{sk}$: $\theta = \pi/2$
    \quad ($m_z = 0$ contour)
\item $r \gg R_{sk}$: $\theta \to 0$
    \quad (background, $m_z = +1$)
\end{itemize}

The skyrmion radius $R_{sk}$ is not a free parameter --- it is
determined by the balance of exchange, DMI, and anisotropy energies.
In MuMax3, this equilibrium is found numerically: an arbitrary
initial guess is evolved under the LLG equation with high damping
until the system relaxes to the energy minimum.  Our simulator
performs the same relaxation during the $J = 0$ phase.  The
profile~\eqref{eq:sk_profile} simply provides a physically
motivated initial condition that is already close to equilibrium,
minimizing the relaxation time needed.

%======================================================================
\section{The Constraint $\mb{m}\cdot d\mb{m}/dt = 0$}
\label{app:m_perp_mdot}
%======================================================================

The identity $\mb{m}\cdot d\mb{m}/dt = 0$ is used throughout
Section~\ref{sec:llgs_explicit_derivation} (and is the geometric
reason the LLGS equation preserves $|\mb{m}|=1$). It follows from two
steps.

\subsection*{Step 1: algebraic identity (always true)}

By the product rule applied to $|\mb{m}|^2 = \mb{m}\cdot\mb{m}$:
%
\begin{equation}
\frac{d|\mb{m}|^2}{dt}
    = \frac{d(\mb{m}\cdot\mb{m})}{dt}
    = \frac{d\mb{m}}{dt}\cdot\mb{m}
        + \mb{m}\cdot\frac{d\mb{m}}{dt}
    = 2\,\mb{m}\cdot\frac{d\mb{m}}{dt}.
\end{equation}
%
Hence
%
\begin{equation}\label{eq:m_dot_mdot_general}
\mb{m}\cdot\frac{d\mb{m}}{dt}
    = \frac{1}{2}\frac{d|\mb{m}|^2}{dt},
\end{equation}
%
which holds for \emph{any} vector $\mb{m}(t)$, with or without the
unit-norm constraint.

\subsection*{Step 2: apply the constraint $|\mb{m}| = 1$}

If $|\mb{m}| = 1$ at all times then $|\mb{m}|^2 = 1$ is a constant,
so
%
\begin{equation}
\frac{d|\mb{m}|^2}{dt} = \frac{d(1)}{dt} = 0.
\end{equation}
%
Substituting into~\eqref{eq:m_dot_mdot_general}:
%
\begin{equation}\label{eq:m_perp_mdot}
\boxed{\;
\mb{m}\cdot\frac{d\mb{m}}{dt} = 0.
\;}
\end{equation}

\subsection*{Geometric meaning}

The state vector $\mb{m}$ lives on the unit two-sphere
$S^2 \subset \mathbb{R}^3$. The tangent space at any point
$\mb{m} \in S^2$ is the plane perpendicular to the radial direction
$\mb{m}$ itself. So any velocity $d\mb{m}/dt$ that keeps $\mb{m}$ on
$S^2$ must be perpendicular to $\mb{m}$, which is precisely
statement~\eqref{eq:m_perp_mdot}.

\begin{center}
\begin{tabular}{ccc}
\textbf{Sphere $S^2$} & $\;$ & \textbf{Tangent plane at $\mb{m}$} \\
                     &      &                                 \\
$|\mb{m}| = 1$       &      & $d\mb{m}/dt \perp \mb{m}$       \\
\end{tabular}
\end{center}

\subsection*{Why the LLGS RHS preserves the constraint}

Every term in the LLGS equation~\eqref{eq:llgs_explicit} has the form
$\mb{m}\times(\,\cdots\,)$ --- precession, damping, and both SOT
torques are cross products with $\mb{m}$. A cross product
$\mb{m}\times\mb{X}$ is orthogonal to both factors, so for any vector
$\mb{X}$:
%
\begin{equation}
\mb{m}\cdot(\mb{m}\times\mb{X}) = 0.
\end{equation}
%
Therefore dotting the entire LLGS RHS with $\mb{m}$ gives zero,
$\mb{m}\cdot d\mb{m}/dt = 0$ is enforced by the equation of motion,
and the unit-norm constraint $|\mb{m}|=1$ is conserved
\emph{exactly} by the continuous-time dynamics. The RK4 integrator
introduces a small $\mathcal{O}(\Delta t^5)$ drift that is removed
by an explicit renormalization step
$\mb{m} \leftarrow \mb{m}/|\mb{m}|$ at the end of each step (see
\texttt{normalize()} in \texttt{src/simulator/integrator.py}).

%======================================================================
\begin{thebibliography}{1}
\bibitem{pham2024}
V.~T.~Pham, N.~Sisodia, I.~Di~Manici,
J.~Urrestarazu-Larra\~{n}aga et al.,
``Fast current-induced skyrmion motion in synthetic
antiferromagnets,''
\emph{Nature Physics} (2024).

\bibitem{mallinson1987}
J.~C.~Mallinson,
``On damped gyromagnetic precession,''
\emph{IEEE Transactions on Magnetics}
\textbf{23}, 2003--2004 (1987).

\bibitem{khvalkovskiy2013}
A.~V.~Khvalkovskiy, V.~Cros, D.~Apalkov, V.~Nikitin,
M.~Krounbi, K.~A.~Zvezdin, A.~Anane, J.~Grollier, A.~Fert,
``Matching domain-wall configuration and spin-orbit torques for
efficient domain-wall motion,''
\emph{Physical Review B} \textbf{87}, 020402(R) (2013).
\end{thebibliography}

\end{document}
